% ============================================================
%  Section 4 — Methodology  (~2 pages)
% ============================================================
\section{Methodology}

\subsection{The QCanvas Pipeline}

\begin{figure}[t]
\centering
\begin{tikzpicture}[
  node distance=0.6cm and 1.1cm,
  box/.style={draw, rounded corners=4pt, minimum width=2.0cm,
              minimum height=0.75cm, align=center, font=\small},
  arr/.style={-{Stealth[length=6pt]}, thick}
]
  \node[box, fill=blue!10]   (q)  {Qiskit\\Circuit};
  \node[box, fill=green!10,  below=of q] (c)  {Cirq\\Circuit};
  \node[box, fill=orange!10, below=of c] (p)  {PennyLane\\QNode};
  \node[box, fill=gray!15,   right=1.3cm of c] (qc) {QCanvas\\Converter};
  \node[box, fill=purple!10, right=1.3cm of qc] (oa) {OpenQASM 3.0\\IR};
  \node[box, fill=red!10,    right=1.3cm of oa] (qs) {QSim\\Backend};
  \node[box, fill=yellow!10, right=1.3cm of qs] (m)  {Metrics\\Engine};
  \draw[arr] (q) -- (qc);
  \draw[arr] (c) -- (qc);
  \draw[arr] (p) -- (qc);
  \draw[arr] (qc) -- (oa);
  \draw[arr] (oa) -- (qs);
  \draw[arr] (qs) -- (m);
\end{tikzpicture}
\caption{The QCanvas unified benchmarking pipeline.
  All three framework circuits are compiled to a common OpenQASM 3.0
  intermediate representation before execution on the shared QSim
  statevector backend.
  Every measured difference is attributable to the framework's
  transpilation strategy.}
\label{fig:pipeline}
\end{figure}

Figure~\ref{fig:pipeline} shows the end-to-end pipeline.
A practitioner writes a circuit idiomatically in their chosen framework.
QCanvas's framework-specific converter traverses the circuit's abstract
syntax tree and emits semantically equivalent OpenQASM 3.0.
The resulting QASM string is passed to QSim for statevector simulation.
Because the intermediate representation and the backend are fixed across
all three frameworks, any difference in compilation latency, gate count,
depth, fidelity, or output distribution is attributable solely to the
framework's transpilation choices.

\subsection{Software Environment and Versions}

Table~\ref{tab:versions} lists the exact package versions used in all
experiments.
Quantum SDKs evolve rapidly and compiler behaviour can change between
minor releases; pinning these versions is a prerequisite for
reproducibility.
The full conda/pip environment file is included in the repository
(\texttt{environment.yml}).

\begin{table}[t]
\caption{Exact software versions used in all benchmarks.}
\label{tab:versions}
\centering
\begin{tabular}{lll}
\toprule
\textbf{Package}          & \textbf{Version} & \textbf{Role} \\
\midrule
Qiskit                    & 1.0.2            & Framework / transpiler \\
Cirq                      & 1.3.0            & Framework / compiler \\
PennyLane                 & 0.35.1           & Framework / decomposer \\
QSim (qsimcirq)           & 0.20.0           & Shared simulation backend \\
Python                    & 3.10.14          & Runtime \\
\bottomrule
\end{tabular}
\end{table}

\subsection{Transpiler Configuration}

Reproducible gate counts require explicit transpiler settings for each
framework.
For \textbf{Qiskit}, all circuits are transpiled at
\texttt{optimization\_level=3} (the maximum level, enabling full
combination of algebraic simplification, commutation analysis, and
Solovay--Kitaev rewrite rules) with no hardware coupling map enforced,
corresponding to an ideal all-to-all connectivity.
For \textbf{Cirq}, circuits are passed through
\texttt{cirq.optimize\_for\_target\_gateset} targeting the
\texttt{CZTargetGateset} with \texttt{allow\_partial\_czs=False};
no additional routing pass is applied (all-to-all topology assumed).
For \textbf{PennyLane}, the \texttt{QNode} is executed on device
\texttt{default.qubit}; the circuit is extracted via \texttt{qml.draw}
after the default decomposition pipeline with no explicit optimisation
pass beyond PennyLane's built-in gate decomposition.

\subsection{Target Gate Basis}

Fair gate-count comparison requires a common target gate vocabulary.
All three frameworks export circuits through the QCanvas converter to
OpenQASM 3.0; the converter normalises every circuit to the basis
$\{\texttt{U3},\,\texttt{CX}\}$---a universal gate set consisting of a
single general-purpose single-qubit rotation and the two-qubit
CNOT~\cite{nielsen2010quantum}.
This normalisation is performed by the \texttt{QASM3Builder} pass inside
QCanvas after framework-specific conversion, ensuring that gate counts
are computed over an identical vocabulary for all three frameworks.
Without this normalisation, a framework emitting \texttt{RX, RY, RZ, CZ}
would be counted differently from one emitting \texttt{U3, CX} even for
an otherwise equivalent circuit.

\subsection{Algorithm Suite}

Table~\ref{tab:algorithms} lists the 12 algorithms in the benchmark
suite.
Each algorithm is implemented \emph{idiomatically} in each framework by
a practitioner familiar with that framework's conventions; no algorithm
is mechanically ported from one framework to another.
Critically, all three implementations encode the \emph{same mathematical
ansatz}---the identical unitary in terms of qubit connectivity and
parameter values before framework-specific decomposition.
The idiomatic coding style governs API calls and gate naming conventions,
not the underlying circuit structure; differences in post-compilation
gate count and depth therefore reflect the framework's compiler and
decomposition choices, not different algorithmic starting points.

\begin{table}[t]
\caption{Algorithm test suite.}
\label{tab:algorithms}
\centering
\begin{tabular}{llc}
\toprule
\textbf{Algorithm}       & \textbf{Category}    & \textbf{Qubit Range} \\
\midrule
Bell State               & Foundational         & 2q \\
GHZ State                & Foundational         & 3--8q \\
Quantum Teleportation    & Communication        & 3q \\
Deutsch--Jozsa           & Oracle               & 3--8q \\
Bernstein--Vazirani      & Oracle               & 3--8q \\
Grover's Algorithm       & Search               & 2--15q (gate counts); 2--6q (simulation) \\
QRNG                     & Randomness           & 4--8q \\
VQE                      & Variational          & 2--6q \\
QAOA                     & Variational          & 2--6q \\
QML XOR Classifier       & Quantum ML           & 2--4q \\
Bit-Flip Error Code      & Error Correction     & 3--9q \\
Phase-Flip Error Code    & Error Correction     & 3--9q \\
\bottomrule
\end{tabular}
\end{table}

\subsection{Metric Categories}

Six metric categories are evaluated for every algorithm--framework pair:
\begin{enumerate}
  \item \textbf{Compilation:} mean compilation latency (ms) $\pm$ std
        over 10 repeated runs; success flag.
  \item \textbf{Structural:} total gate count, circuit depth, multi-qubit
        gate count, multi-qubit ratio, per-gate-type counts.
  \item \textbf{Scaling:} power-law exponent $a$, coefficient $c$,
        coefficient of determination $R^2$, scaling class
        (constant / sub-linear / linear / super-linear).
  \item \textbf{Simulation:} mean and std of simulation time (ms) and
        memory (MB), mean CPU utilisation (\%), fidelity under a
        depolarizing noise model~\cite{nielsen2010quantum}.
  \item \textbf{Statistical equivalence:} pairwise
        JSD~\cite{lin1991divergence} at 4096 shots; equivalence declared
        for JSD $< 0.02$; per-algorithm \texttt{all\_equivalent} verdict.
  \item \textbf{Holistic:} Quantum Volume~\cite{cross2019validating};
        QPack sub-scores $S_{\text{compile\_speed}}$,
        $S_{\text{scalability}}$, $S_{\text{accuracy}}$,
        $S_{\text{capacity}}$, $S_{\text{overall}}$.
\end{enumerate}

\subsection{Measurement Protocol}

All timed measurements are repeated 10 times on the same physical
machine with no concurrent workloads; we report the mean and standard
deviation.
Simulation shot counts of 1024, 4096, and 8192 are evaluated; the
primary analysis uses 4096 shots.
All package versions are pinned (Table~\ref{tab:versions}); the full
conda environment is reproducible from the repository
(\texttt{environment.yml}).
Grover's Algorithm gate counts are measured through compilation only
(without statevector simulation) for qubit ranges 7--15, where full
statevector simulation would exceed available memory ($2^{15}$ amplitudes
require $\approx$256 GB RAM); this allows the power-law scaling curve
to be established across 14 data points rather than 5.
